\chapter{Introduction}\label{ch:intro}

\section{Motivation}\label{sec:motivation}

An old photograph is rarely damaged in just one way. Take a print shot
through a soft lens on grainy film, scanned at low resolution and saved as a
heavily compressed JPEG. It carries all four defects, and each one hits a
different part of the spectrum. Optical softening removes high spatial
frequencies. Resampling throws away detail above a resolution limit and adds
interpolation artefacts. Sensor and film noise spreads energy roughly evenly
across frequencies. In natural images the true signal is weakest at high
frequencies, so that is where noise dominates. JPEG quantises the mid and
high frequency coefficients of each $8\times8$ block, which leaves blocking
and ringing.

Most learned restoration networks treat this mixture as a single problem.
DnCNN \cite{zhang2017dncnn}, a deep convolutional network, sees the whole
degraded image and learns to undo every corruption in it without being told
which is which. Transformer restorers such as Restormer
\cite{zamir2022restormer} work very well, but they have tens of millions of
parameters. An archive, a library or a hobbyist with modest hardware cannot
run them.

We started from a simple idea. If each corruption sits in its own frequency
band, we could split the image into bands explicitly and let a small network
pick a filter for every band at every location. A large kernel suits smooth,
slowly varying structure. A small kernel responds to fine detail and to the
high-frequency part of the noise. A learned gate chooses between them pixel
by pixel. We call the result \fcsg{}, a frequency-conditioned, spatially
gated network. We hold it to a strict budget of fewer than five million
parameters and fewer than ten GFLOPs for a $256\times256$ input.

\section{Problem statement}\label{sec:problem}

Let $f \in [0,1]^{H\times W\times3}$ be a clean RGB image and let
\begin{equation}
	x = \mathcal{D}(f;\, \theta), \qquad
	\theta = (\sigma_\text{blur},\, s,\, \sigma_n,\, q)
	\label{eq:degradation}
\end{equation}
be a degraded observation. $\mathcal{D}$ applies Gaussian blur of width
$\sigma_\text{blur}$, bicubic down- and up-sampling by a factor $s$, additive
Gaussian noise with standard deviation $\sigma_n$ and JPEG compression at
quality $q$. We draw $\theta$ at random for every image and never give it to
the restoration model. The task is to learn a mapping $g_\phi$ whose output
$\hat f = g_\phi(x)$ is close to $f$ in PSNR, SSIM and LPIPS. The network
must also meet two hard limits:
\begin{itemize}
	\item fewer than $5\times10^{6}$ trainable parameters, and
	\item fewer than 10 reported GFLOPs for one forward pass on a
	$256\times256$ RGB input.
\end{itemize}

\section{Hypothesis and objectives}\label{sec:objectives}

We test one hypothesis. \emph{Apply a fixed radial frequency decomposition,
then route each band per pixel to the top two of several shared
convolutional experts with different kernel sizes. The learned routing will
change with image content and with the degradation, the trained network will
depend on it for output quality, and the whole model will stay inside the
budget.} We wrote this as a claim to check, not as a promise that the network
beats anything.

The work has four objectives.
\begin{enumerate}
	\item[\textbf{O1}] Build \fcsg{} and check that it behaves as designed.
	The checks cover exact band reconstruction, the parameter and FLOP
	limits, memory, gradient flow and overfitting a single image
	(Chapters~\ref{ch:method} and~\ref{ch:setup}).
	\item[\textbf{O2}] Get baselines we can trust. We reproduce a published
	FFDNet result and train DnCNN on a Gaussian task with a known outcome
	(Chapter~\ref{ch:setup}).
	\item[\textbf{O3}] Train \fcsg{}, DnCNN and two FFDNet variants under one
	matched composite degradation protocol. Compare quality, size, compute and
	measured latency (Chapter~\ref{ch:results}).
	\item[\textbf{O4}] Describe what the router learned. Where does it send
	each band, how does that relate to the sampled degradation, and how much
	does the trained model rely on it (Chapter~\ref{ch:routing})?
\end{enumerate}

\section{Contributions}\label{sec:contributions}

This report makes four concrete contributions.
\begin{itemize}
	\item A restoration network that routes frequency bands spatially, with
	97,088 parameters and 9.39 reported GFLOPs at $256\times256$. We also give
	the tests showing the code matches the equations.
	\item A matched comparison. We train DnCNN, blind FFDNet, oracle FFDNet and
	\fcsg{} from scratch on the same synthetic composite degradation, with the
	same schedule, the same number of updates and the same evaluation
	protocol.
	\item A routing analysis. We relate per-pixel expert weights to the
	recorded degradation parameters and measure what happens when we override
	the gate at inference.
	\item A plain account of how it turned out. \fcsg{} reaches 27.24\dB{},
	against 27.83\dB{} for DnCNN and 28.19\dB{} for blind FFDNet, with 5.7
	times fewer parameters than DnCNN. It is not faster on a GPU. Its low FLOP
	count did not become low latency.
\end{itemize}
We make no claim of state-of-the-art quality. Frequency decomposition and
expert routing both have a long history (Chapter~\ref{ch:literature}). We
test one specific combination of them under a tight budget.

\section{Organisation of the report}\label{sec:organisation}

Chapter~\ref{ch:literature} reviews restoration networks, frequency-domain
methods and mixture-of-experts routing, and places \fcsg{} among them.
Chapter~\ref{ch:method} describes the degradation model and then the network
one module at a time. Chapter~\ref{ch:setup} covers the data pipeline,
baselines, training and evaluation protocols, and the verification runs.
Chapter~\ref{ch:results} reports restoration quality and computational cost.
Chapter~\ref{ch:routing} looks at the learned routing.
Chapter~\ref{ch:conclusion} sums up the findings and their limits and lists
the experiments we would run next. Appendix~\ref{app:repro} records what
anyone needs to reproduce the work, and Appendix~\ref{app:per-image} lists
results for every image.
